% =====================================================================
%  BUỔI 12 — TÀI LIỆU BUỔI HỌC (từ hoc-sinh/02-tai-lieu-buoi-hoc.md)
%  Sinh tự động bằng scripts/md2tex-chung.py — không sửa tay.
% =====================================================================
\documentclass[11pt, a4paper]{article}
% =====================================================================
%  CAE SHANGHAI — GÓI LỆNH DÙNG CHUNG CHO TÀI LIỆU LUYỆN THI CSCA TOÁN
%  Biên dịch bằng XeLaTeX. Mọi file .tex trong buổi học \input file này.
% =====================================================================

% ---------- NGÔN NGỮ & PHÔNG CHỮ ----------
\usepackage{fontspec}
\usepackage{polyglossia}
\setmainlanguage{vietnamese}
\setotherlanguage{english}

% Chữ Latin (Anh + Việt): TeX Gyre Termes — bản clone Times New Roman,
% đi kèm TeX Live, hỗ trợ đầy đủ dấu tiếng Việt.
\setmainfont{TeX Gyre Termes}
\setsansfont{TeX Gyre Heros}

% Chữ Hán: ưu tiên Source Han Serif SC (思源宋体) — chuẩn học thuật Trung Quốc.
% Nếu máy chưa cài, tự động dùng Songti SC (宋体, có sẵn trên macOS).
% Cả hai đều thuộc họ 宋体 nên diện mạo tài liệu không đổi.
\IfFontExistsTF{Source Han Serif SC}{
  \newcommand{\cjkfamilyname}{Source Han Serif SC}
}{
  \IfFontExistsTF{Songti SC}{
    \newcommand{\cjkfamilyname}{Songti SC}
  }{
    \newcommand{\cjkfamilyname}{Noto Serif CJK SC}
  }
}

% xeCJK tự chuyển font khi gặp ký tự Hán, không cần bọc \cjk{} thủ công.
% Đây là điều kiện bắt buộc: font Latin (TeX Gyre Termes) không có glyph chữ Hán,
% nếu không khai báo xeCJK thì mọi chữ Hán viết trần sẽ bị nuốt khỏi bản in.
\usepackage{xeCJK}
\setCJKfamilyfont{caehan}{\cjkfamilyname}[Scale=0.95]
\setCJKmainfont{\cjkfamilyname}[Scale=0.95]
\setCJKsansfont{\cjkfamilyname}[Scale=0.95]
\setCJKmonofont{\cjkfamilyname}[Scale=0.95]

% \cjk{} ép dùng font Hán trong ngữ cảnh cần thiết (ví dụ trong tiêu đề in đậm
% hoặc trong bảng có định dạng riêng). Dùng \CJKfamily chứ không phải tên font,
% vì tên font chỉ là văn bản — sẽ in ra chữ thay vì đổi font.
\newcommand{\cjk}[1]{{\CJKfamily{caehan}#1}}

% ---------- TOÁN HỌC ----------
\usepackage{amsmath, amssymb, amsthm, mathtools}
\usepackage{unicode-math}
\setmathfont{TeX Gyre Termes Math}

% ---------- BỐ CỤC TRANG ----------
% Lề phải chừa vùng an toàn của letterhead:
%   banner trên chiếm 0–19mm  → top tối thiểu 2,2cm
%   banner dưới chiếm 281–297mm → bottom tối thiểu 2,2cm
% Lề trên 2,6cm (26mm) để tiêu đề đầu trang không đè banner (banner hết ở 19mm).
\usepackage[a4paper, top=2.6cm, bottom=2.4cm, left=2.5cm, right=2.0cm]{geometry}
\usepackage{fancyhdr}
\usepackage{titlesec}
\usepackage{enumitem}
\usepackage{graphicx}
\usepackage{booktabs, array, longtable, tabularx}
\usepackage{xcolor}
\usepackage[most]{tcolorbox}
\usepackage[colorlinks=true, linkcolor=cae_blue, citecolor=cae_blue, urlcolor=cae_blue]{hyperref}

% ---------- BẢNG MÀU THƯƠNG HIỆU ----------
\definecolor{cae_blue}{RGB}{0, 51, 102}
\definecolor{cae_gold}{RGB}{196, 149, 60}
\definecolor{box_bg}{RGB}{245, 247, 250}
\definecolor{box_line}{RGB}{0, 51, 102}

% ---------- TIÊU ĐỀ MỤC ----------
% Đệm đầu trang: banner letterhead chiếm 0–19mm, mà lề trên là 22mm.
% \section* đầu trang (sau \clearpage) sát lề nên phải chừa thêm, nếu không
% sẽ đè lên banner.
\titleformat{\section}
  {\normalfont\large\bfseries\sffamily\color{cae_blue}}{\thesection.}{0.6em}{}
\titleformat{\subsection}
  {\normalfont\normalsize\bfseries\sffamily\color{cae_blue!85!black}}{\thesubsection}{0.6em}{}
\titlespacing*{\section}{0pt}{14pt}{6pt}
\titlespacing*{\subsection}{0pt}{10pt}{4pt}

% Tiêu đề không đánh số ở đầu trang: \section* sau \clearpage nằm sát lề trên
% nên đè lên banner letterhead. Dùng \vspace* ngay TRƯỚC \section* không đủ,
% vì titlesec chèn khoảng cách riêng. Cách chắc chắn là đệm bằng \null.
\newcommand{\sectiondautrang}[1]{\section*{#1}}

% Dạng bài dùng \subsection*{} nên cần định dạng riêng, đậm và có gạch chân màu
\titleformat{\subsubsection}
  {\normalfont\normalsize\bfseries\sffamily\color{cae_blue!85!black}}{\thesubsubsection}{0.6em}{}

% ---------- LETTERHEAD (ảnh nền toàn trang) ----------
% Ảnh letterhead của CAE Shanghai: banner xanh trên/dưới kèm logo chìm giữa.
% Tỉ lệ ảnh 1414x2000 = 0.7070, khớp A4 (0.7071) nên chèn thẳng không méo.
% Đặt ở lớp nền dưới cùng để chữ không bị ảnh che.
% Dùng \AddToShipoutPictureBG (không có dấu *) để áp dụng cho MỌI trang.
\usepackage{eso-pic}
\AddToShipoutPictureBG{%
  \AtPageLowerLeft{%
    \includegraphics[width=\paperwidth,height=\paperheight]{images/letterhead.png}%
  }%
}

% ---------- ĐẦU / CHÂN TRANG ----------
% Letterhead đã có banner trên và dưới, nên bỏ dòng chạy để tránh chồng chữ.
% Số trang đặt trong vùng nội dung, sát trên banner dưới, lề phải.
\pagestyle{fancy}
\setlength{\headheight}{0pt}
\fancyhf{}
\renewcommand{\headrulewidth}{0pt}
\renewcommand{\footrulewidth}{0pt}
\fancyfoot[R]{\raisebox{0.9cm}{\small\color{cae_blue}\thepage}}

% ---------- MÔI TRƯỜNG NỘI DUNG ----------
% Trong tcolorbox, LaTeX đặt parindent = parskip = 0, nên các đoạn văn trong
% khung dính liền nhau. Đặt lại khoảng cách đoạn cho mọi khung.
\newcommand{\khoangcachdoan}{\setlength{\parskip}{0.5em}\setlength{\parindent}{0pt}}
\tcbset{before upper={\khoangcachdoan}}

\newtcolorbox{dinhnghia}[1][]{
  colback=box_bg, colframe=box_line, boxrule=0.6pt, arc=1pt,
  fonttitle=\bfseries, title={Định nghĩa #1}, breakable}
\newtcolorbox{dinhly}[1][]{
  colback=box_bg, colframe=cae_gold, boxrule=0.6pt, arc=1pt,
  fonttitle=\bfseries, title={Tính chất #1}, breakable}
\newtcolorbox{vidu}[1][]{
  colback=white, colframe=black!40, boxrule=0.4pt, arc=1pt,
  fonttitle=\bfseries, title={Ví dụ #1}, breakable}
\newtcolorbox{baitap}[1][]{
  colback=white, colframe=cae_blue!50, boxrule=0.4pt, arc=1pt,
  fonttitle=\bfseries, title={Câu #1}, breakable}

% Ô đáp án: viền xanh đậm hơn, dùng trong file đáp án
\newtcolorbox{dapan}[1][]{
  colback=box_bg, colframe=cae_blue, boxrule=0.6pt, arc=1pt,
  fonttitle=\bfseries, title={Câu #1}, breakable}

% ---------- TÊN CHỨNG MINH ----------
\renewcommand{\proofname}{\normalfont\bfseries Chứng minh}

% ---------- MÔN HỌC ----------
% Đổi được ở file .tex trước khi % =====================================================================
%  CAE SHANGHAI — GÓI LỆNH DÙNG CHUNG CHO TÀI LIỆU LUYỆN THI CSCA TOÁN
%  Biên dịch bằng XeLaTeX. Mọi file .tex trong buổi học \input file này.
% =====================================================================

% ---------- NGÔN NGỮ & PHÔNG CHỮ ----------
\usepackage{fontspec}
\usepackage{polyglossia}
\setmainlanguage{vietnamese}
\setotherlanguage{english}

% Chữ Latin (Anh + Việt): TeX Gyre Termes — bản clone Times New Roman,
% đi kèm TeX Live, hỗ trợ đầy đủ dấu tiếng Việt.
\setmainfont{TeX Gyre Termes}
\setsansfont{TeX Gyre Heros}

% Chữ Hán: ưu tiên Source Han Serif SC (思源宋体) — chuẩn học thuật Trung Quốc.
% Nếu máy chưa cài, tự động dùng Songti SC (宋体, có sẵn trên macOS).
% Cả hai đều thuộc họ 宋体 nên diện mạo tài liệu không đổi.
\IfFontExistsTF{Source Han Serif SC}{
  \newcommand{\cjkfamilyname}{Source Han Serif SC}
}{
  \IfFontExistsTF{Songti SC}{
    \newcommand{\cjkfamilyname}{Songti SC}
  }{
    \newcommand{\cjkfamilyname}{Noto Serif CJK SC}
  }
}

% xeCJK tự chuyển font khi gặp ký tự Hán, không cần bọc \cjk{} thủ công.
% Đây là điều kiện bắt buộc: font Latin (TeX Gyre Termes) không có glyph chữ Hán,
% nếu không khai báo xeCJK thì mọi chữ Hán viết trần sẽ bị nuốt khỏi bản in.
\usepackage{xeCJK}
\setCJKfamilyfont{caehan}{\cjkfamilyname}[Scale=0.95]
\setCJKmainfont{\cjkfamilyname}[Scale=0.95]
\setCJKsansfont{\cjkfamilyname}[Scale=0.95]
\setCJKmonofont{\cjkfamilyname}[Scale=0.95]

% \cjk{} ép dùng font Hán trong ngữ cảnh cần thiết (ví dụ trong tiêu đề in đậm
% hoặc trong bảng có định dạng riêng). Dùng \CJKfamily chứ không phải tên font,
% vì tên font chỉ là văn bản — sẽ in ra chữ thay vì đổi font.
\newcommand{\cjk}[1]{{\CJKfamily{caehan}#1}}

% ---------- TOÁN HỌC ----------
\usepackage{amsmath, amssymb, amsthm, mathtools}
\usepackage{unicode-math}
\setmathfont{TeX Gyre Termes Math}

% ---------- BỐ CỤC TRANG ----------
% Lề phải chừa vùng an toàn của letterhead:
%   banner trên chiếm 0–19mm  → top tối thiểu 2,2cm
%   banner dưới chiếm 281–297mm → bottom tối thiểu 2,2cm
% Lề trên 2,6cm (26mm) để tiêu đề đầu trang không đè banner (banner hết ở 19mm).
\usepackage[a4paper, top=2.6cm, bottom=2.4cm, left=2.5cm, right=2.0cm]{geometry}
\usepackage{fancyhdr}
\usepackage{titlesec}
\usepackage{enumitem}
\usepackage{graphicx}
\usepackage{booktabs, array, longtable, tabularx}
\usepackage{xcolor}
\usepackage[most]{tcolorbox}
\usepackage[colorlinks=true, linkcolor=cae_blue, citecolor=cae_blue, urlcolor=cae_blue]{hyperref}

% ---------- BẢNG MÀU THƯƠNG HIỆU ----------
\definecolor{cae_blue}{RGB}{0, 51, 102}
\definecolor{cae_gold}{RGB}{196, 149, 60}
\definecolor{box_bg}{RGB}{245, 247, 250}
\definecolor{box_line}{RGB}{0, 51, 102}

% ---------- TIÊU ĐỀ MỤC ----------
% Đệm đầu trang: banner letterhead chiếm 0–19mm, mà lề trên là 22mm.
% \section* đầu trang (sau \clearpage) sát lề nên phải chừa thêm, nếu không
% sẽ đè lên banner.
\titleformat{\section}
  {\normalfont\large\bfseries\sffamily\color{cae_blue}}{\thesection.}{0.6em}{}
\titleformat{\subsection}
  {\normalfont\normalsize\bfseries\sffamily\color{cae_blue!85!black}}{\thesubsection}{0.6em}{}
\titlespacing*{\section}{0pt}{14pt}{6pt}
\titlespacing*{\subsection}{0pt}{10pt}{4pt}

% Tiêu đề không đánh số ở đầu trang: \section* sau \clearpage nằm sát lề trên
% nên đè lên banner letterhead. Dùng \vspace* ngay TRƯỚC \section* không đủ,
% vì titlesec chèn khoảng cách riêng. Cách chắc chắn là đệm bằng \null.
\newcommand{\sectiondautrang}[1]{\section*{#1}}

% Dạng bài dùng \subsection*{} nên cần định dạng riêng, đậm và có gạch chân màu
\titleformat{\subsubsection}
  {\normalfont\normalsize\bfseries\sffamily\color{cae_blue!85!black}}{\thesubsubsection}{0.6em}{}

% ---------- LETTERHEAD (ảnh nền toàn trang) ----------
% Ảnh letterhead của CAE Shanghai: banner xanh trên/dưới kèm logo chìm giữa.
% Tỉ lệ ảnh 1414x2000 = 0.7070, khớp A4 (0.7071) nên chèn thẳng không méo.
% Đặt ở lớp nền dưới cùng để chữ không bị ảnh che.
% Dùng \AddToShipoutPictureBG (không có dấu *) để áp dụng cho MỌI trang.
\usepackage{eso-pic}
\AddToShipoutPictureBG{%
  \AtPageLowerLeft{%
    \includegraphics[width=\paperwidth,height=\paperheight]{images/letterhead.png}%
  }%
}

% ---------- ĐẦU / CHÂN TRANG ----------
% Letterhead đã có banner trên và dưới, nên bỏ dòng chạy để tránh chồng chữ.
% Số trang đặt trong vùng nội dung, sát trên banner dưới, lề phải.
\pagestyle{fancy}
\setlength{\headheight}{0pt}
\fancyhf{}
\renewcommand{\headrulewidth}{0pt}
\renewcommand{\footrulewidth}{0pt}
\fancyfoot[R]{\raisebox{0.9cm}{\small\color{cae_blue}\thepage}}

% ---------- MÔI TRƯỜNG NỘI DUNG ----------
% Trong tcolorbox, LaTeX đặt parindent = parskip = 0, nên các đoạn văn trong
% khung dính liền nhau. Đặt lại khoảng cách đoạn cho mọi khung.
\newcommand{\khoangcachdoan}{\setlength{\parskip}{0.5em}\setlength{\parindent}{0pt}}
\tcbset{before upper={\khoangcachdoan}}

\newtcolorbox{dinhnghia}[1][]{
  colback=box_bg, colframe=box_line, boxrule=0.6pt, arc=1pt,
  fonttitle=\bfseries, title={Định nghĩa #1}, breakable}
\newtcolorbox{dinhly}[1][]{
  colback=box_bg, colframe=cae_gold, boxrule=0.6pt, arc=1pt,
  fonttitle=\bfseries, title={Tính chất #1}, breakable}
\newtcolorbox{vidu}[1][]{
  colback=white, colframe=black!40, boxrule=0.4pt, arc=1pt,
  fonttitle=\bfseries, title={Ví dụ #1}, breakable}
\newtcolorbox{baitap}[1][]{
  colback=white, colframe=cae_blue!50, boxrule=0.4pt, arc=1pt,
  fonttitle=\bfseries, title={Câu #1}, breakable}

% Ô đáp án: viền xanh đậm hơn, dùng trong file đáp án
\newtcolorbox{dapan}[1][]{
  colback=box_bg, colframe=cae_blue, boxrule=0.6pt, arc=1pt,
  fonttitle=\bfseries, title={Câu #1}, breakable}

% ---------- TÊN CHỨNG MINH ----------
\renewcommand{\proofname}{\normalfont\bfseries Chứng minh}

% ---------- MÔN HỌC ----------
% Đổi được ở file .tex trước khi % =====================================================================
%  CAE SHANGHAI — GÓI LỆNH DÙNG CHUNG CHO TÀI LIỆU LUYỆN THI CSCA TOÁN
%  Biên dịch bằng XeLaTeX. Mọi file .tex trong buổi học \input file này.
% =====================================================================

% ---------- NGÔN NGỮ & PHÔNG CHỮ ----------
\usepackage{fontspec}
\usepackage{polyglossia}
\setmainlanguage{vietnamese}
\setotherlanguage{english}

% Chữ Latin (Anh + Việt): TeX Gyre Termes — bản clone Times New Roman,
% đi kèm TeX Live, hỗ trợ đầy đủ dấu tiếng Việt.
\setmainfont{TeX Gyre Termes}
\setsansfont{TeX Gyre Heros}

% Chữ Hán: ưu tiên Source Han Serif SC (思源宋体) — chuẩn học thuật Trung Quốc.
% Nếu máy chưa cài, tự động dùng Songti SC (宋体, có sẵn trên macOS).
% Cả hai đều thuộc họ 宋体 nên diện mạo tài liệu không đổi.
\IfFontExistsTF{Source Han Serif SC}{
  \newcommand{\cjkfamilyname}{Source Han Serif SC}
}{
  \IfFontExistsTF{Songti SC}{
    \newcommand{\cjkfamilyname}{Songti SC}
  }{
    \newcommand{\cjkfamilyname}{Noto Serif CJK SC}
  }
}

% xeCJK tự chuyển font khi gặp ký tự Hán, không cần bọc \cjk{} thủ công.
% Đây là điều kiện bắt buộc: font Latin (TeX Gyre Termes) không có glyph chữ Hán,
% nếu không khai báo xeCJK thì mọi chữ Hán viết trần sẽ bị nuốt khỏi bản in.
\usepackage{xeCJK}
\setCJKfamilyfont{caehan}{\cjkfamilyname}[Scale=0.95]
\setCJKmainfont{\cjkfamilyname}[Scale=0.95]
\setCJKsansfont{\cjkfamilyname}[Scale=0.95]
\setCJKmonofont{\cjkfamilyname}[Scale=0.95]

% \cjk{} ép dùng font Hán trong ngữ cảnh cần thiết (ví dụ trong tiêu đề in đậm
% hoặc trong bảng có định dạng riêng). Dùng \CJKfamily chứ không phải tên font,
% vì tên font chỉ là văn bản — sẽ in ra chữ thay vì đổi font.
\newcommand{\cjk}[1]{{\CJKfamily{caehan}#1}}

% ---------- TOÁN HỌC ----------
\usepackage{amsmath, amssymb, amsthm, mathtools}
\usepackage{unicode-math}
\setmathfont{TeX Gyre Termes Math}

% ---------- BỐ CỤC TRANG ----------
% Lề phải chừa vùng an toàn của letterhead:
%   banner trên chiếm 0–19mm  → top tối thiểu 2,2cm
%   banner dưới chiếm 281–297mm → bottom tối thiểu 2,2cm
% Lề trên 2,6cm (26mm) để tiêu đề đầu trang không đè banner (banner hết ở 19mm).
\usepackage[a4paper, top=2.6cm, bottom=2.4cm, left=2.5cm, right=2.0cm]{geometry}
\usepackage{fancyhdr}
\usepackage{titlesec}
\usepackage{enumitem}
\usepackage{graphicx}
\usepackage{booktabs, array, longtable, tabularx}
\usepackage{xcolor}
\usepackage[most]{tcolorbox}
\usepackage[colorlinks=true, linkcolor=cae_blue, citecolor=cae_blue, urlcolor=cae_blue]{hyperref}

% ---------- BẢNG MÀU THƯƠNG HIỆU ----------
\definecolor{cae_blue}{RGB}{0, 51, 102}
\definecolor{cae_gold}{RGB}{196, 149, 60}
\definecolor{box_bg}{RGB}{245, 247, 250}
\definecolor{box_line}{RGB}{0, 51, 102}

% ---------- TIÊU ĐỀ MỤC ----------
% Đệm đầu trang: banner letterhead chiếm 0–19mm, mà lề trên là 22mm.
% \section* đầu trang (sau \clearpage) sát lề nên phải chừa thêm, nếu không
% sẽ đè lên banner.
\titleformat{\section}
  {\normalfont\large\bfseries\sffamily\color{cae_blue}}{\thesection.}{0.6em}{}
\titleformat{\subsection}
  {\normalfont\normalsize\bfseries\sffamily\color{cae_blue!85!black}}{\thesubsection}{0.6em}{}
\titlespacing*{\section}{0pt}{14pt}{6pt}
\titlespacing*{\subsection}{0pt}{10pt}{4pt}

% Tiêu đề không đánh số ở đầu trang: \section* sau \clearpage nằm sát lề trên
% nên đè lên banner letterhead. Dùng \vspace* ngay TRƯỚC \section* không đủ,
% vì titlesec chèn khoảng cách riêng. Cách chắc chắn là đệm bằng \null.
\newcommand{\sectiondautrang}[1]{\section*{#1}}

% Dạng bài dùng \subsection*{} nên cần định dạng riêng, đậm và có gạch chân màu
\titleformat{\subsubsection}
  {\normalfont\normalsize\bfseries\sffamily\color{cae_blue!85!black}}{\thesubsubsection}{0.6em}{}

% ---------- LETTERHEAD (ảnh nền toàn trang) ----------
% Ảnh letterhead của CAE Shanghai: banner xanh trên/dưới kèm logo chìm giữa.
% Tỉ lệ ảnh 1414x2000 = 0.7070, khớp A4 (0.7071) nên chèn thẳng không méo.
% Đặt ở lớp nền dưới cùng để chữ không bị ảnh che.
% Dùng \AddToShipoutPictureBG (không có dấu *) để áp dụng cho MỌI trang.
\usepackage{eso-pic}
\AddToShipoutPictureBG{%
  \AtPageLowerLeft{%
    \includegraphics[width=\paperwidth,height=\paperheight]{images/letterhead.png}%
  }%
}

% ---------- ĐẦU / CHÂN TRANG ----------
% Letterhead đã có banner trên và dưới, nên bỏ dòng chạy để tránh chồng chữ.
% Số trang đặt trong vùng nội dung, sát trên banner dưới, lề phải.
\pagestyle{fancy}
\setlength{\headheight}{0pt}
\fancyhf{}
\renewcommand{\headrulewidth}{0pt}
\renewcommand{\footrulewidth}{0pt}
\fancyfoot[R]{\raisebox{0.9cm}{\small\color{cae_blue}\thepage}}

% ---------- MÔI TRƯỜNG NỘI DUNG ----------
% Trong tcolorbox, LaTeX đặt parindent = parskip = 0, nên các đoạn văn trong
% khung dính liền nhau. Đặt lại khoảng cách đoạn cho mọi khung.
\newcommand{\khoangcachdoan}{\setlength{\parskip}{0.5em}\setlength{\parindent}{0pt}}
\tcbset{before upper={\khoangcachdoan}}

\newtcolorbox{dinhnghia}[1][]{
  colback=box_bg, colframe=box_line, boxrule=0.6pt, arc=1pt,
  fonttitle=\bfseries, title={Định nghĩa #1}, breakable}
\newtcolorbox{dinhly}[1][]{
  colback=box_bg, colframe=cae_gold, boxrule=0.6pt, arc=1pt,
  fonttitle=\bfseries, title={Tính chất #1}, breakable}
\newtcolorbox{vidu}[1][]{
  colback=white, colframe=black!40, boxrule=0.4pt, arc=1pt,
  fonttitle=\bfseries, title={Ví dụ #1}, breakable}
\newtcolorbox{baitap}[1][]{
  colback=white, colframe=cae_blue!50, boxrule=0.4pt, arc=1pt,
  fonttitle=\bfseries, title={Câu #1}, breakable}

% Ô đáp án: viền xanh đậm hơn, dùng trong file đáp án
\newtcolorbox{dapan}[1][]{
  colback=box_bg, colframe=cae_blue, boxrule=0.6pt, arc=1pt,
  fonttitle=\bfseries, title={Câu #1}, breakable}

% ---------- TÊN CHỨNG MINH ----------
\renewcommand{\proofname}{\normalfont\bfseries Chứng minh}

% ---------- MÔN HỌC ----------
% Đổi được ở file .tex trước khi \input{cae-style} bằng \renewcommand{\monhoc}{...}
\providecommand{\monhoc}{Môn Toán}

% ---------- TRANG BÌA GỌN ----------
% #1 = loại tài liệu (TÀI LIỆU CHUẨN BỊ, TÀI LIỆU BUỔI HỌC, ...)
% #2 = số buổi
% #3 = tên buổi tiếng Việt
% #4 = tên buổi tiếng Hán
% #5 = module và tỉ trọng
\newcommand{\trangbia}[5]{%
  \begin{center}
    \vspace*{0.6cm}
    {\large\bfseries\color{cae_blue} CAE SHANGHAI}\\[0.15cm]
    {\footnotesize Tài liệu luyện thi CSCA · \monhoc}\\[0.5cm]
    {\color{cae_gold}\rule{0.55\linewidth}{0.8pt}}\\[0.35cm]
    {\small\bfseries\color{cae_blue!80!black} #1}\\[0.35cm]
    {\Large\bfseries\color{cae_blue} BUỔI #2 — #3}\\[0.2cm]
    {\large\color{cae_blue!85!black} \cjk{#4}}\\[0.25cm]
    \ifx\relax#5\relax\else{\small #5}\\[0.35cm]\fi
    {\color{cae_gold}\rule{0.55\linewidth}{0.8pt}}\\[0.3cm]
    {\footnotesize Biên soạn: Hoàng Dương, CAE Shanghai \quad·\quad \today}
  \end{center}
  \vspace{0.3cm}
}

% ---------- TRANG BÌA CHO ĐỀ KIỂM TRA ----------
% Khác \trangbia ở chỗ không có chữ "BUỔI N" — dùng cho đề thi, không phải bài giảng.
% #1 = loại tài liệu, #2 = tên kỳ thi (Việt), #3 = tên kỳ thi (Hán), #4 = thông số
\newcommand{\biade}[4]{%
  \begin{center}
    \vspace*{0.6cm}
    {\large\bfseries\color{cae_blue} CAE SHANGHAI}\\[0.15cm]
    {\footnotesize Tài liệu luyện thi CSCA · \monhoc}\\[0.5cm]
    {\color{cae_gold}\rule{0.55\linewidth}{0.8pt}}\\[0.35cm]
    {\small\bfseries\color{cae_blue!80!black} #1}\\[0.35cm]
    {\Large\bfseries\color{cae_blue} #2}\\[0.2cm]
    {\large\color{cae_blue!85!black} \cjk{#3}}\\[0.25cm]
    {\small #4}\\[0.35cm]
    {\color{cae_gold}\rule{0.55\linewidth}{0.8pt}}\\[0.3cm]
    {\footnotesize Biên soạn: Hoàng Dương, CAE Shanghai \quad·\quad \today}
  \end{center}
  \vspace{0.3cm}
}

% ---------- TIỆN ÍCH ----------
% Nhãn từ khóa in đậm ở đầu đoạn, dùng cho Từ vựng / Từ khóa nhận dạng / Chú ý
% Khoảng cách đoạn do khung tự đặt (xem \khoangcachdoan ở trên),
% ở đây chỉ cần tách đoạn và bỏ thụt đầu dòng.
\newcommand{\nhan}[1]{\par\noindent\textbf{#1.}\ \ignorespaces}

% Bảng gọn, không đánh số
\newcommand{\bang}[1]{%
  \begin{center}#1\end{center}}

% Đánh dấu ô trống cho học sinh điền
\newcommand{\chotrong}{\makebox[3.2cm]{\dotfill}}
 bằng \renewcommand{\monhoc}{...}
\providecommand{\monhoc}{Môn Toán}

% ---------- TRANG BÌA GỌN ----------
% #1 = loại tài liệu (TÀI LIỆU CHUẨN BỊ, TÀI LIỆU BUỔI HỌC, ...)
% #2 = số buổi
% #3 = tên buổi tiếng Việt
% #4 = tên buổi tiếng Hán
% #5 = module và tỉ trọng
\newcommand{\trangbia}[5]{%
  \begin{center}
    \vspace*{0.6cm}
    {\large\bfseries\color{cae_blue} CAE SHANGHAI}\\[0.15cm]
    {\footnotesize Tài liệu luyện thi CSCA · \monhoc}\\[0.5cm]
    {\color{cae_gold}\rule{0.55\linewidth}{0.8pt}}\\[0.35cm]
    {\small\bfseries\color{cae_blue!80!black} #1}\\[0.35cm]
    {\Large\bfseries\color{cae_blue} BUỔI #2 — #3}\\[0.2cm]
    {\large\color{cae_blue!85!black} \cjk{#4}}\\[0.25cm]
    \ifx\relax#5\relax\else{\small #5}\\[0.35cm]\fi
    {\color{cae_gold}\rule{0.55\linewidth}{0.8pt}}\\[0.3cm]
    {\footnotesize Biên soạn: Hoàng Dương, CAE Shanghai \quad·\quad \today}
  \end{center}
  \vspace{0.3cm}
}

% ---------- TRANG BÌA CHO ĐỀ KIỂM TRA ----------
% Khác \trangbia ở chỗ không có chữ "BUỔI N" — dùng cho đề thi, không phải bài giảng.
% #1 = loại tài liệu, #2 = tên kỳ thi (Việt), #3 = tên kỳ thi (Hán), #4 = thông số
\newcommand{\biade}[4]{%
  \begin{center}
    \vspace*{0.6cm}
    {\large\bfseries\color{cae_blue} CAE SHANGHAI}\\[0.15cm]
    {\footnotesize Tài liệu luyện thi CSCA · \monhoc}\\[0.5cm]
    {\color{cae_gold}\rule{0.55\linewidth}{0.8pt}}\\[0.35cm]
    {\small\bfseries\color{cae_blue!80!black} #1}\\[0.35cm]
    {\Large\bfseries\color{cae_blue} #2}\\[0.2cm]
    {\large\color{cae_blue!85!black} \cjk{#3}}\\[0.25cm]
    {\small #4}\\[0.35cm]
    {\color{cae_gold}\rule{0.55\linewidth}{0.8pt}}\\[0.3cm]
    {\footnotesize Biên soạn: Hoàng Dương, CAE Shanghai \quad·\quad \today}
  \end{center}
  \vspace{0.3cm}
}

% ---------- TIỆN ÍCH ----------
% Nhãn từ khóa in đậm ở đầu đoạn, dùng cho Từ vựng / Từ khóa nhận dạng / Chú ý
% Khoảng cách đoạn do khung tự đặt (xem \khoangcachdoan ở trên),
% ở đây chỉ cần tách đoạn và bỏ thụt đầu dòng.
\newcommand{\nhan}[1]{\par\noindent\textbf{#1.}\ \ignorespaces}

% Bảng gọn, không đánh số
\newcommand{\bang}[1]{%
  \begin{center}#1\end{center}}

% Đánh dấu ô trống cho học sinh điền
\newcommand{\chotrong}{\makebox[3.2cm]{\dotfill}}
 bằng \renewcommand{\monhoc}{...}
\providecommand{\monhoc}{Môn Toán}

% ---------- TRANG BÌA GỌN ----------
% #1 = loại tài liệu (TÀI LIỆU CHUẨN BỊ, TÀI LIỆU BUỔI HỌC, ...)
% #2 = số buổi
% #3 = tên buổi tiếng Việt
% #4 = tên buổi tiếng Hán
% #5 = module và tỉ trọng
\newcommand{\trangbia}[5]{%
  \begin{center}
    \vspace*{0.6cm}
    {\large\bfseries\color{cae_blue} CAE SHANGHAI}\\[0.15cm]
    {\footnotesize Tài liệu luyện thi CSCA · \monhoc}\\[0.5cm]
    {\color{cae_gold}\rule{0.55\linewidth}{0.8pt}}\\[0.35cm]
    {\small\bfseries\color{cae_blue!80!black} #1}\\[0.35cm]
    {\Large\bfseries\color{cae_blue} BUỔI #2 — #3}\\[0.2cm]
    {\large\color{cae_blue!85!black} \cjk{#4}}\\[0.25cm]
    \ifx\relax#5\relax\else{\small #5}\\[0.35cm]\fi
    {\color{cae_gold}\rule{0.55\linewidth}{0.8pt}}\\[0.3cm]
    {\footnotesize Biên soạn: Hoàng Dương, CAE Shanghai \quad·\quad \today}
  \end{center}
  \vspace{0.3cm}
}

% ---------- TRANG BÌA CHO ĐỀ KIỂM TRA ----------
% Khác \trangbia ở chỗ không có chữ "BUỔI N" — dùng cho đề thi, không phải bài giảng.
% #1 = loại tài liệu, #2 = tên kỳ thi (Việt), #3 = tên kỳ thi (Hán), #4 = thông số
\newcommand{\biade}[4]{%
  \begin{center}
    \vspace*{0.6cm}
    {\large\bfseries\color{cae_blue} CAE SHANGHAI}\\[0.15cm]
    {\footnotesize Tài liệu luyện thi CSCA · \monhoc}\\[0.5cm]
    {\color{cae_gold}\rule{0.55\linewidth}{0.8pt}}\\[0.35cm]
    {\small\bfseries\color{cae_blue!80!black} #1}\\[0.35cm]
    {\Large\bfseries\color{cae_blue} #2}\\[0.2cm]
    {\large\color{cae_blue!85!black} \cjk{#3}}\\[0.25cm]
    {\small #4}\\[0.35cm]
    {\color{cae_gold}\rule{0.55\linewidth}{0.8pt}}\\[0.3cm]
    {\footnotesize Biên soạn: Hoàng Dương, CAE Shanghai \quad·\quad \today}
  \end{center}
  \vspace{0.3cm}
}

% ---------- TIỆN ÍCH ----------
% Nhãn từ khóa in đậm ở đầu đoạn, dùng cho Từ vựng / Từ khóa nhận dạng / Chú ý
% Khoảng cách đoạn do khung tự đặt (xem \khoangcachdoan ở trên),
% ở đây chỉ cần tách đoạn và bỏ thụt đầu dòng.
\newcommand{\nhan}[1]{\par\noindent\textbf{#1.}\ \ignorespaces}

% Bảng gọn, không đánh số
\newcommand{\bang}[1]{%
  \begin{center}#1\end{center}}

% Đánh dấu ô trống cho học sinh điền
\newcommand{\chotrong}{\makebox[3.2cm]{\dotfill}}


\begin{document}

\trangbia{TÀI LIỆU BUỔI HỌC}{12}{VECTOR VÀ SỐ PHỨC}{向量与复数}{Module M3 \quad·\quad Tỉ trọng đề thi $\sim$6\% \quad·\quad khoảng 2–4 câu}

\section{TỪ KHÓA NHẬN DẠNG ĐỀ}

Đọc được từ khóa là xác định được dạng bài mà không cần dịch cả đề. Bảng dưới đây gom những từ khóa quyết định của chuyên đề này.

\bang{
\begin{tabularx}{\linewidth}{@{\hspace{6pt}}llX@{\hspace{6pt}}}
\toprule
\textbf{Từ khóa} & \textbf{Nghĩa} & \textbf{Dạng bài tương ứng} \\
\midrule
数量积 & tích vô hướng & Tính $a \cdot b$ từ tọa độ \\
夹角 & góc & Tính góc giữa hai vector \\
投影 & hình chiếu & Chiếu vector này lên phương vector kia \\
模 & môđun & Tính độ dài vector hoặc môđun số phức \\
垂直 & vuông góc & Tìm tham số để $a \cdot b = 0$ \\
平行 / 共线 & song song / cùng phương & Tìm tham số để $x_1y_2 - x_2y_1 = 0$ \\
中点 & trung điểm & Tọa độ trung điểm, trung tuyến \\
实部 & phần thực & Đọc phần thực của số phức \\
虚部 & phần ảo & Đọc phần ảo của số phức \\
共轭复数 & số phức liên hợp & Đổi dấu phần ảo \\
纯虚数 & số thuần ảo & Điều kiện: phần thực bằng 0 và phần ảo khác 0 \\
虚数单位 & đơn vị ảo & Rút gọn lũy thừa của $i$ theo chu kì 4 \\
\bottomrule
\end{tabularx}}

\nhan{Chú ý} Ba cặp từ khóa dễ đọc nhầm nhất:

\bang{
\begin{tabularx}{\linewidth}{@{\hspace{6pt}}lX@{\hspace{6pt}}}
\toprule
\textbf{Cặp từ khóa} & \textbf{Khác nhau ở đâu} \\
\midrule
模 và 共轭复数 & 模 là một số thực không âm; 共轭复数 là một số phức \\
纯虚数 và 虚数 & 纯虚数 bắt buộc phần thực bằng $0$; 虚数 chỉ cần phần ảo khác $0$ \\
垂直 và 平行 & 垂直 dùng $a \cdot b = 0$; 平行 dùng $x_1y_2 = x_2y_1$ \\
\bottomrule
\end{tabularx}}

\section{KIẾN THỨC TRỌNG TÂM}

\subsection{Vector trong mặt phẳng}

\begin{dinhnghia}[{1 — Vector và môđun}]
Vector là đại lượng có cả độ lớn và hướng. Vector $a = (x, y)$ có môđun $\lvert a \rvert = \sqrt{x^2 + y^2}$.
\end{dinhnghia}
\begin{dinhnghia}[{2 — Tích vô hướng}]
Với $a = (x_1, y_1)$ và $b = (x_2, y_2)$:

$$a \cdot b = x_1x_2 + y_1y_2 = \lvert a \rvert \lvert b \rvert \cos\theta$$

trong đó $\theta$ là góc giữa hai vector.
\end{dinhnghia}
\begin{dinhly}[{1 — Các phép toán tọa độ}]
Với $a = (x_1, y_1)$, $b = (x_2, y_2)$ và số thực $\lambda$:

$$a + b = (x_1 + x_2,\ y_1 + y_2), \qquad a - b = (x_1 - x_2,\ y_1 - y_2), \qquad \lambda a = (\lambda x_1,\ \lambda y_1)$$
\end{dinhly}
\begin{dinhly}[{2 — Điều kiện vuông góc}]
Với hai vector khác vector không:

$$a \perp b \iff a \cdot b = 0 \iff x_1x_2 + y_1y_2 = 0$$
\end{dinhly}
\begin{dinhly}[{3 — Điều kiện song song}]
Với $b$ khác vector không:

$$a \parallel b \iff x_1y_2 - x_2y_1 = 0$$
\end{dinhly}
\begin{dinhly}[{4 — Công thức góc}]
Góc $\theta$ giữa hai vector khác vector không xác định bởi:

$$\cos\theta = \frac{a \cdot b}{\lvert a \rvert \lvert b \rvert} = \frac{x_1x_2 + y_1y_2}{\sqrt{x_1^2 + y_1^2}\sqrt{x_2^2 + y_2^2}}$$
\end{dinhly}
\begin{dinhly}[{5 — Hình chiếu}]
Hình chiếu vô hướng của $a$ lên phương $b$ bằng:

$$\frac{a \cdot b}{\lvert b \rvert}$$

\nhan{Chú ý} Hình chiếu là một \textbf{số} có thể âm, không phải một vector. Đề dùng chữ 投影 mà hỏi một giá trị trong bốn phương án số thì đó là hình chiếu vô hướng.
\end{dinhly}
\begin{dinhly}[{6 — Tọa độ trung điểm}]
Trung điểm $M$ của đoạn $AB$ với $A(x_1, y_1)$, $B(x_2, y_2)$ có tọa độ:

$$M = \left(\frac{x_1 + x_2}{2},\ \frac{y_1 + y_2}{2}\right)$$

\nhan{Chú ý} Vector $\overrightarrow{AB} = B - A$, tức lấy tọa độ điểm cuối trừ tọa độ điểm đầu. Đảo thứ tự là sai dấu cả hai thành phần, và đây là lỗi mất điểm phổ biến nhất ở dạng tìm tọa độ điểm.

\nhan{Cách nhanh} Hai vector cùng phương khi tỉ số các thành phần tương ứng bằng nhau. Với $a = (1,2)$, $b = (2,4)$: $\dfrac{2}{1} = \dfrac{4}{2} = 2$ nên cùng phương. Cách này nhanh hơn viết phương trình $x_1y_2 = x_2y_1$ nhưng phải kiểm tra trường hợp có thành phần bằng $0$.
\end{dinhly}
\subsection{Số phức}

\begin{dinhnghia}[{3 — Số phức và các thành phần}]
Số phức có dạng $z = a + bi$ với $a, b \in \mathbb{R}$ và $i^2 = -1$. Số $a$ gọi là \textbf{phần thực}, số $b$ gọi là \textbf{phần ảo}.
\end{dinhnghia}
\begin{dinhnghia}[{4 — Số phức liên hợp}]
Số phức liên hợp của $z = a + bi$ là $\overline{z} = a - bi$.
\end{dinhnghia}
\begin{dinhnghia}[{5 — Môđun số phức}]
Môđun của $z = a + bi$ là:

$$\lvert z \rvert = \sqrt{a^2 + b^2}$$
\end{dinhnghia}
\begin{dinhly}[{7 — Bốn phép toán}]
Với $z_1 = a + bi$ và $z_2 = c + di$:

$$z_1 \pm z_2 = (a \pm c) + (b \pm d)i$$

$$z_1 \cdot z_2 = (ac - bd) + (ad + bc)i$$

$$\frac{z_1}{z_2} = \frac{(a + bi)(c - di)}{c^2 + d^2} \quad (z_2 \ne 0)$$
\end{dinhly}
\begin{dinhly}[{8 — Chu kì lũy thừa của $i$}]
Với mọi số nguyên $n$:

$$i^{4n} = 1, \qquad i^{4n+1} = i, \qquad i^{4n+2} = -1, \qquad i^{4n+3} = -i$$

\nhan{Cách nhanh} Chia số mũ cho $4$ và chỉ quan tâm số dư. Với $i^{2025}$: $2025 = 4 \times 506 + 1$ nên $i^{2025} = i$. Không cần viết chu kì ra giấy.
\end{dinhly}
\begin{dinhly}[{9 — Các đẳng thức về môđun và liên hợp}]
$$z \cdot \overline{z} = \lvert z \rvert^2 = a^2 + b^2, \qquad \lvert \overline{z} \rvert = \lvert z \rvert, \qquad \overline{z_1 \cdot z_2} = \overline{z_1} \cdot \overline{z_2}$$
\end{dinhly}
\begin{dinhly}[{10 — Điều kiện số thực và số thuần ảo}]
Cho $z = a + bi$ với $a, b \in \mathbb{R}$:

$$z \in \mathbb{R} \iff b = 0; \qquad z \ \text{là số thuần ảo} \iff a = 0 \ \text{và} \ b \ne 0$$

\nhan{Bẫy} Số thuần ảo yêu cầu \textbf{hai} điều kiện. Nếu chỉ đặt $a = 0$ mà quên $b \ne 0$, ta nhận luôn cả $z = 0$, mà số $0$ là số thực chứ không phải số thuần ảo.

\nhan{Chú ý} Phần ảo của $z = a + bi$ là số thực $b$, \textbf{không phải} $bi$. Đề hỏi 虚部 thì đáp án là một số thực; phương án dạng $4i$ là phương án nhiễu.

\nhan{Cách nhanh} Khi đề cho $z = \dfrac{a + bi}{c + di}$, đừng khai triển ngay. Nhân tử và mẫu với $c - di$ rồi tính $c^2 + d^2$ trước để biết mẫu số, sau đó mới rút gọn tử.

\nhan{Nhận xét} Phép chia số phức và phép tính môđun là hai dạng chiếm phần lớn số câu về số phức trong đề thi thật. Cả hai đều chỉ dùng một công thức duy nhất, không đòi hỏi biến đổi sáng tạo.

\end{dinhly}
\section{CÁC DẠNG BÀI}

\subsection{Dạng 1 — Tích vô hướng và góc giữa hai vector}

\nhan{Cách nhận dạng} Đề cho tọa độ hai vector và hỏi 数量积 (tích vô hướng), 夹角 (góc), hoặc 余弦值 (giá trị cosin).

\nhan{Cách làm} Tính $a \cdot b = x_1x_2 + y_1y_2$, tính hai môđun, thay vào $\cos\theta = \dfrac{a \cdot b}{\lvert a \rvert \lvert b \rvert}$. Nếu đề cho độ dài và góc thì đi ngược lại: dùng $a \cdot b = \lvert a \rvert \lvert b \rvert \cos\theta$.

\begin{vidu}[{1}]
Cho $u = (1, \sqrt{3})$ và $v = (\sqrt{3}, 1)$. Tìm góc giữa hai vector.

Tích vô hướng: $u \cdot v = 1 \cdot \sqrt{3} + \sqrt{3} \cdot 1 = 2\sqrt{3}$.

Hai môđun: $\lvert u \rvert = \sqrt{1 + 3} = 2$ và $\lvert v \rvert = \sqrt{3 + 1} = 2$.

$$\cos\theta = \frac{2\sqrt{3}}{2 \cdot 2} = \frac{\sqrt{3}}{2}$$

Vậy $\theta = 30^\circ$.

\end{vidu}
\subsection{Dạng 2 — Tìm tham số để hai vector vuông góc}

\nhan{Cách nhận dạng} Đề có tham số $k$, $m$, $t$ và cụm 垂直 (vuông góc).

\nhan{Cách làm} Lập phương trình $x_1x_2 + y_1y_2 = 0$, giải phương trình bậc nhất hoặc bậc hai theo tham số. Chỉ cần một bước, không vẽ hình.

\begin{vidu}[{2}]
Cho $a = (5, 1)$ và $b = (-2, 2k)$. Tìm $k$ để $a \perp b$.

Điều kiện vuông góc:

$$5 \cdot (-2) + 1 \cdot 2k = 0 \iff -10 + 2k = 0 \iff k = 5$$

\end{vidu}
\subsection{Dạng 3 — Tìm tham số để hai vector song song}

\nhan{Cách nhận dạng} Đề có cụm 平行 (song song) hoặc 共线 (cùng phương).

\nhan{Cách làm} Lập phương trình $x_1y_2 - x_2y_1 = 0$. Với vector hai chiều, phương trình này là bậc nhất theo tham số và cho một nghiệm duy nhất. Với vector ba chiều thì xét tỉ số ba thành phần.

\begin{vidu}[{3}]
Cho $a = (k, 1)$ và $b = (4, k)$. Tìm $k$ để $a \parallel b$.

Điều kiện song song:

$$k \cdot k - 4 \cdot 1 = 0 \iff k^2 = 4 \iff k = \pm 2$$

\nhan{Bẫy} Đáp án chỉ ghi $k = 2$ là thiếu nghiệm âm. Phương trình bậc hai theo tham số luôn cho hai nghiệm, phải viết đủ.

\end{vidu}
\subsection{Dạng 4 — Môđun, hình chiếu và tọa độ trung điểm}

\nhan{Cách nhận dạng} Đề hỏi 模 (môđun), 投影 (hình chiếu), 中点 (trung điểm), hoặc cho tọa độ điểm đầu và vector rồi hỏi tọa độ điểm cuối.

\nhan{Cách làm} Môđun dùng căn của tổng bình phương. Hình chiếu của $a$ lên $b$ bằng $\dfrac{a \cdot b}{\lvert b \rvert}$. Tọa độ điểm cuối bằng tọa độ điểm đầu cộng vector.

\begin{vidu}[{4}]
Cho $\overrightarrow{AB} = (3, -6)$ và $A(-1, 2)$. Tìm tọa độ $B$.

Vì $\overrightarrow{AB} = B - A$ nên $B = A + \overrightarrow{AB}$:

$$B = (-1 + 3,\ 2 + (-6)) = (2, -4)$$

\nhan{Chú ý} Công thức này viết dưới dạng tọa độ: điểm cuối bằng điểm đầu cộng vector. Nếu đề cho điểm cuối và hỏi điểm đầu thì làm phép trừ.

\end{vidu}
\subsection{Dạng 5 — Phép toán cộng, trừ, nhân, chia số phức}

\nhan{Cách nhận dạng} Đề cho hai hoặc ba số phức và yêu cầu tính $z_1 \pm z_2$, $z_1 z_2$, hoặc $\dfrac{z_1}{z_2}$.

\nhan{Cách làm} Cộng trừ theo từng phần. Nhân như nhân đa thức rồi thay $i^2 = -1$. Chia bằng cách nhân tử và mẫu với số phức liên hợp của mẫu.

\begin{vidu}[{5}]
Cho $z_1 = 4 - i$ và $z_2 = (2 + 5i)i$. Tính $z_1 - z_2$.

Trước hết rút gọn $z_2$:

$$z_2 = (2 + 5i)i = 2i + 5i^2 = 2i - 5 = -5 + 2i$$

$$z_1 - z_2 = (4 - i) - (-5 + 2i) = 9 - 3i$$

\nhan{Cách nhanh} Luôn rút gọn số phức có chứa $i$ nhân bên ngoài trước khi thực hiện phép toán chính. Bỏ bước này là phải nhân hai lần.

\end{vidu}
\subsection{Dạng 6 — Phần thực, phần ảo, số phức liên hợp}

\nhan{Cách nhận dạng} Đề hỏi 实部, 虚部, 共轭复数, hoặc cho một biểu thức và hỏi phần ảo của kết quả.

\nhan{Cách làm} Đưa số phức về dạng chuẩn $a + bi$ bằng cách thực hiện hết các phép toán, sau đó đọc trực tiếp $a$ và $b$.

\begin{vidu}[{6}]
Cho $z = \dfrac{2+i}{1-i}$. Tìm số phức liên hợp $\overline{z}$.

Nhân tử và mẫu với $1 + i$:

$$z = \frac{(2+i)(1+i)}{(1-i)(1+i)} = \frac{2 + 2i + i + i^2}{1 - i^2} = \frac{1 + 3i}{2} = \frac{1}{2} + \frac{3}{2}i$$

$$\overline{z} = \frac{1}{2} - \frac{3}{2}i$$

\end{vidu}
\subsection{Dạng 7 — Tham số để số phức là số thực hoặc số thuần ảo}

\nhan{Cách nhận dạng} Đề cho $z$ chứa tham số $a$, $m$, $k$ và cụm 纯虚数 (thuần ảo) hoặc 实数 (số thực).

\nhan{Cách làm} Đưa $z$ về dạng $a + bi$ với $a$ và $b$ là biểu thức chứa tham số, rồi áp dụng Tính chất 10. Số thuần ảo cần hai điều kiện, số thực chỉ cần một.

\begin{vidu}[{7}]
Cho $z = (a + i)(1 - i)$. Tìm $a$ để $z$ là số thuần ảo.

Khai triển:

$$z = (a + i)(1 - i) = a - ai + i - i^2 = (a + 1) + (1 - a)i$$

Phần thực $a + 1$, phần ảo $1 - a$. Điều kiện số thuần ảo:

$$a + 1 = 0 \implies a = -1$$

Kiểm tra điều kiện thứ hai: $1 - (-1) = 2 \ne 0$ nên $a = -1$ thỏa mãn.

\end{vidu}
\subsection{Dạng 8 — Môđun số phức và lũy thừa của đơn vị ảo}

\nhan{Cách nhận dạng} Đề hỏi 模, $\lvert z \rvert$, $\lvert \overline{z} \rvert$, $z \cdot \overline{z}$, hoặc yêu cầu tính $i^n$ với $n$ lớn.

\nhan{Cách làm} Với môđun, đưa $z$ về dạng $a + bi$ rồi tính $\sqrt{a^2 + b^2}$. Với $i^n$, chia $n$ cho $4$ và dùng số dư.

\begin{vidu}[{8}]
Tính môđun của $z = -3 - 4i$.

$$\lvert z \rvert = \sqrt{(-3)^2 + (-4)^2} = \sqrt{9 + 16} = 5$$

\nhan{Cách nhanh} Bộ số $3$–$4$–$5$ và $5$–$12$–$13$ xuất hiện rất nhiều trong đề. Nhìn thấy phần thực và phần ảo là bội của $3$ và $4$ thì đọc ngay đáp án mà không cần bấm căn.

\nhan{Chú ý} $\lvert \overline{z} \rvert = \lvert z \rvert$, nên đề hỏi môđun của số phức liên hợp thì tính môđun của $z$ gốc, không cần đổi dấu.

\end{vidu}
\section{PHỤ LỤC — CẤU TRÚC ĐỀ THI CSCA MÔN TOÁN}

\bang{
\begin{tabularx}{\linewidth}{@{\hspace{6pt}}lX@{\hspace{6pt}}}
\toprule
\textbf{Đặc điểm} & \textbf{Thông số} \\
\midrule
Số câu & 48 câu trắc nghiệm \\
Thời gian & 60 phút \\
Thang điểm & 100 điểm \\
Ngôn ngữ & Tiếng Trung hoặc tiếng Anh (thí sinh chọn) \\
Máy tính & Không được dùng \\
\bottomrule
\end{tabularx}}

Thời gian trung bình mỗi câu là \textbf{1 phút 15 giây}. Chuyên đề Vector và số phức chỉ chiếm 2–4 câu, nhưng đây là nhóm câu \textbf{dễ lấy điểm nhất} trong toàn đề: hầu hết chỉ cần một công thức và một phép thay số. Mục tiêu là xử lý mỗi câu trong khoảng 40 giây.
\vspace{0.8cm}
\begin{center}
{\footnotesize\color{cae_blue} Tài liệu nội bộ, CAE SHANGHAI}
\end{center}

\end{document}
